\section{Consejos para la generación joven}

\subsection{La elección de dirección de investigación en 2025}

Para quienes cursan la universidad o comienzan un doctorado en robótica hoy, la recomendación de Yang Shuo (杨硕) es:

\begin{knowledgebox}{Las direcciones centrales de la investigación en robótica en 2025}
\begin{enumerate}
    \item \textbf{Los fundamentos siguen siendo importantes}: aún hace falta estudiar sistemas y señales, pero no es necesario agotar, como hacían las generaciones anteriores, todas las derivaciones de cada método tradicional; basta con entender su conexión con el conocimiento matemático fundamental
    \item \textbf{La dirección más importante}: dentro del marco de las redes neuronales, cómo extraer estructuras de baja dimensión a partir de datos de sensores, datos de imagen, etc., para ayudar al robot a ejecutar mejor sus tareas; esta dirección general nunca cambiará
    \item \textbf{No aferrarse a la forma humanoide}: lo que el usuario necesita puede ser un martillo sencillo y no una navaja suiza compleja; hay que partir del problema y no de la tecnología
\end{enumerate}
\end{knowledgebox}

\subsection{Sentido de misión y disfrute}

Al final de la entrevista, el presentador le preguntó a Yang Shuo cuál era la reward function de su vida. La respuesta de Yang Shuo se divide en etapas:
\begin{itemize}
    \item En DJI: hacer un muy buen producto
    \item Durante el doctorado: hacer un buen trabajo de investigación
    \item Después de tener hijos: el bienestar de la familia se volvió lo más importante
    \item En la etapa actual: primero, cuidar bien a la familia y a los hijos; segundo, explorar, a través de Miaodong Technology (妙动科技), los productos robóticos y el problema de la simbiosis entre humanos y máquinas
\end{itemize}

Él considera que «al llegar a cierta edad, en un mismo periodo solo se pueden hacer bien dos o tres cosas».

Sobre el origen de su motivación, Yang Shuo dijo: «Soy alguien arrastrado sin remedio por esta tendencia, a ayudar a que los robots crezcan». Su sentido de misión tiene tres dimensiones: ayudar a que los robots crezcan, ayudar a los amigos humanos a entender cómo convivir con los robots, y ayudar a los robots del futuro a aprender a convivir con los humanos.

Al final, el presentador comentó con emoción: «Desde tu primer artículo en Zhihu hasta el último, no he sentido que la persona detrás del texto haya cambiado mucho. Tu visión de la sociedad y la tecnología en su conjunto nunca ha vacilado».

\begin{warningbox}{El árbol de habilidades puede trastocarse, pero la intención original no}
La síntesis del presentador transmite un mensaje central: las direcciones tecnológicas y el árbol de habilidades que en 2025 nos parecen muy importantes podrían quedar por completo trastocados dentro de diez años. Pero lo verdaderamente importante, la voluntad de «querer resolver este problema», no se trastoca; mientras consideremos importante la robótica, esa perseverancia nos guiará a través de los cambios de paradigma tecnológico.
\end{warningbox}

\subsection{Resumen de la sección}

El mensaje central que Yang Shuo deja a la generación joven puede resumirse así: mantener una base sólida sin necesidad de agotar los detalles menores, enfocarse en la dirección general invariable de «extraer estructuras de baja dimensión a partir de los datos para apoyar la ejecución de tareas», generar disfrute a partir del sentido de misión, y profundizar de manera constante en aquello a lo que uno se dedica. Las rutas tecnológicas pueden cambiar, pero la pasión por los problemas y la determinación de resolverlos no quedarán obsoletas.
